\subsection*{CU-03: Jugar como invitado}
\addcontentsline{toc}{subsection}{CU-03: Jugar como invitado}
\label{cu-03}

\begin{fichacu}{CU-03}{Jugar como invitado}
  \crudFila{Responsable}{José Eduardo Prior Hernández}
  \crudFila{Actor principal}{Invitado}
  \crudFila{Actores de sistema}{Servidor de partidas, Base de datos}
  \crudFila{Prototipos}{6c}
  \crudFila{Prioridad}{Must (debe estar en la entrega). Categoría (a) Obligatorio: característica 4 del cliente (jugar como invitado).}
  \crudFila{Objetivo}{Permitir que una persona empiece a jugar sin dar correo ni contraseña, creando una cuenta de tipo invitado con la que puede crear salas y unirse a ellas, jugar en ellas y chatear. La cuenta de invitado no busca partida clasificatoria y no aparece en el ranking.}
  \crudFila{Disparador}{El Invitado selecciona ``Jugar como invitado'' en la \texttt{GUI\_\allowbreak{}Login}.}
  \textbf{Precondiciones} & \cuCelda{\begin{cureglas}
      \item \textbf{\mbox{PRE-01}} El cliente del SJM está instalado y en ejecución.
      \item \textbf{\mbox{PRE-02}} El Servidor de partidas está en ejecución y tiene conexión disponible con la Base de datos.
      \item \textbf{\mbox{PRE-03}} El cliente no conserva una cuenta de invitado anterior en este dispositivo. \textit{(Ver \mbox{FA-01})}
    \end{cureglas}} \\ \hline
  \textbf{Flujo normal} & \cuCelda{\begin{cuenum}[start=1]
      \item El Invitado selecciona ``Jugar como invitado'' en la \texttt{GUI\_\allowbreak{}Login}.
      \item El SJM despliega la \texttt{GUI\_\allowbreak{}MessageConfirm} advirtiendo que la cuenta de invitado vive solo en este dispositivo y que sus partidas no cuentan para el ranking, con las opciones ``Continuar'' y ``Cancelar'' (\mbox{RN-01}). \textit{(Ver \mbox{FA-02})}
      \item El Invitado selecciona ``Continuar''.
      \item El SJM abre la conexión TCP con el Servidor de partidas y negocia la versión del protocolo. \textit{(Ver \mbox{EX-01}, \mbox{EX-02})}
      \item El SJM envía el mensaje \texttt{GUEST\_\allowbreak{}SESSION\_\allowbreak{}REQUEST} con la huella del dispositivo y la versión del cliente. \textit{(Ver \mbox{EX-05}, \mbox{EX-06})}
      \item El Servidor de partidas comprueba cuántas cuentas de invitado se han creado desde esa dirección de red en la última hora (\mbox{RN-07}). \textit{(Ver \mbox{FA-03})}
      \item El Servidor de partidas genera un \texttt{nickname} provisional único, con el formato de invitado y un sufijo numérico (\mbox{RN-02}). \textit{(Ver \mbox{FA-04})}
      \item El Servidor de partidas inicia una transacción sobre la Base de datos.
    \end{cuenum}} \\
   & \cuCelda{\begin{cuenum}[start=9]
      \item El Servidor de partidas crea el registro de \textit{Usuario} con \texttt{tipo\_\allowbreak{}cuenta} en \texttt{INVITADO}, \texttt{estado\_\allowbreak{}cuenta} en \texttt{ACTIVA}, el \texttt{nickname} provisional, la \texttt{fecha\_\allowbreak{}registro} actual, y \texttt{correo}, \texttt{contrasena\_\allowbreak{}hash} y \texttt{contrasena\_\allowbreak{}sal} en nulo (\mbox{RN-03}). No crea \textit{EstadisticaModo} (\mbox{RN-06}). \textit{(Ver \mbox{EX-03})}
      \item El Servidor de partidas crea la \textit{Sesion}, con su token, la \texttt{fecha\_\allowbreak{}inicio}, la \texttt{fecha\_\allowbreak{}expiracion}, la dirección de red de origen, la huella y el nombre del dispositivo y la versión del cliente. \textit{(Ver \mbox{EX-03})}
      \item El Servidor de partidas confirma la transacción. \textit{(Ver \mbox{EX-04})}
      \item El Servidor de partidas registra el alta en la \textit{BitacoraAcceso} con el resultado \texttt{ALTA\_\allowbreak{}INVITADO}.
      \item El Servidor de partidas responde con \texttt{GUEST\_\allowbreak{}SESSION\_\allowbreak{}OK}, que contiene el token de sesión y el \texttt{nickname} provisional.
      \item El SJM guarda el token en el dispositivo, de modo que al volver a abrir el juego se recupere la misma cuenta de invitado y no se cree otra (\mbox{RN-04}).
      \item El SJM despliega la \texttt{GUI\_\allowbreak{}MainMenu} con una marca visible de cuenta de invitado (\mbox{RN-01}).
      \item Fin del caso de uso.
    \end{cuenum}} \\ \hline
  \textbf{Flujos alternos} & \cuCelda{\textbf{\mbox{FA-01}: El dispositivo ya tiene una cuenta de invitado}\par
    \begin{cuenum}[start=1]
      \item El SJM envía \texttt{GUEST\_\allowbreak{}SESSION\_\allowbreak{}RESUME} con el token guardado en lugar de solicitar una cuenta nueva.
      \item El Servidor de partidas recupera el \textit{Usuario} de tipo \texttt{INVITADO} asociado y verifica que siga en \texttt{ACTIVA}. \textit{(Ver \mbox{FA-05}, \mbox{FA-06})}
      \item El Servidor de partidas crea una \textit{Sesion} nueva para ese \textit{Usuario}.
      \item El flujo continúa en el paso 12 del FN. El Invitado recupera su \texttt{nickname} provisional y su historial de partidas: es la misma cuenta, no una nueva.
    \end{cuenum}} \\
   & \cuCelda{\textbf{\mbox{FA-02}: El Invitado cancela el aviso}\par
    \begin{cuenum}[start=1]
      \item El Invitado selecciona ``Cancelar''.
      \item El SJM vuelve a la \texttt{GUI\_\allowbreak{}Login} sin enviar nada al Servidor de partidas.
      \item Fin del caso de uso.
    \end{cuenum}} \\
   & \cuCelda{\textbf{\mbox{FA-03}: Se alcanzó el límite de altas de invitado desde esa dirección}\par
    \begin{cuenum}[start=1]
      \item El Servidor de partidas responde con \texttt{GUEST\_\allowbreak{}LIMIT\_\allowbreak{}REACHED} y los minutos que faltan.
      \item El SJM muestra la \texttt{GUI\_\allowbreak{}MessageWarning} indicando que se alcanzó el límite y proponiendo crear una cuenta, que no tiene límite.
      \item Si el Invitado acepta la propuesta, el flujo continúa en el \mbox{CU-02} Registrar cuenta.
      \item Fin del caso de uso.
    \end{cuenum}} \\
   & \cuCelda{\textbf{\mbox{FA-04}: El nickname provisional generado ya está en uso}\par
    \begin{cuenum}[start=1]
      \item El Servidor de partidas genera otro sufijo y vuelve a comprobar.
      \item Si tras varios intentos no consigue uno libre, responde con \texttt{SERVER\_\allowbreak{}BUSY} y el flujo continúa en el paso 1 del FN.
      \item En caso contrario, el flujo continúa en el paso 8 del FN.
    \end{cuenum}} \\
   & \cuCelda{\textbf{\mbox{FA-05}: La cuenta de invitado guardada ya no sirve}\par
    \begin{cuenum}[start=1]
      \item El Servidor de partidas encuentra que el \textit{Usuario} de tipo \texttt{INVITADO} fue eliminado, por ejemplo al caducar (\mbox{RN-09}), o que tiene una sanción de ámbito de cuenta vigente.
      \item El Servidor de partidas responde con \texttt{GUEST\_\allowbreak{}ACCOUNT\_\allowbreak{}UNAVAILABLE} y el motivo.
      \item El SJM descarta el token guardado y el flujo continúa en el paso 2 del FN, creando una cuenta de invitado nueva. Si el motivo fue una sanción, en cambio, el SJM conserva el token, despliega la \texttt{GUI\_\allowbreak{}BannedAccount} con el motivo, el número de strike y cuándo termina, igual que el \mbox{FA-09} del \mbox{CU-01}, y termina el caso de uso.
    \end{cuenum}} \\
   & \cuCelda{\textbf{\mbox{FA-06}: La suspensión de la cuenta de invitado ya venció}\par
    \begin{cuenum}[start=1]
      \item El Servidor de partidas encuentra el \texttt{estado\_\allowbreak{}cuenta} en \texttt{SUSPENDIDA} y ninguna \textit{Sancion} de ámbito \texttt{CUENTA} con \texttt{activa} verdadero y \texttt{fecha\_\allowbreak{}fin} posterior a la fecha actual.
      \item El Servidor de partidas devuelve el \texttt{estado\_\allowbreak{}cuenta} a \texttt{ACTIVA} y registra el cambio en la \textit{BitacoraAcceso} con el resultado \texttt{SUSPENSION\_\allowbreak{}VENCIDA}, igual que el \mbox{RN-16} del \mbox{CU-01} para las cuentas registradas (\mbox{RN-10}). \textit{(Ver \mbox{EX-03})}
      \item El flujo continúa en el paso 3 del \mbox{FA-01}. Si la sanción sigue vigente, el flujo es el del \mbox{FA-05}.
    \end{cuenum}} \\ \hline
  \textbf{Excepciones} & \cuCelda{\textbf{\mbox{EX-01}: No se puede establecer la conexión con el servidor}\par
    \begin{cuenum}[start=1]
      \item El SJM detecta que la conexión TCP no pudo establecerse dentro del tiempo límite y registra el fallo en la bitácora local.
      \item El SJM muestra la \texttt{GUI\_\allowbreak{}MessageError} con las opciones ``Reintentar'' y ``Aceptar''.
      \item Si el Invitado selecciona ``Reintentar'', el flujo continúa en el paso 4 del FN. Si selecciona ``Aceptar'', continúa en el paso 1.
    \end{cuenum}} \\
   & \cuCelda{\textbf{\mbox{EX-02}: La versión del protocolo es incompatible}\par
    \begin{cuenum}[start=1]
      \item El Servidor de partidas responde con \texttt{PROTOCOL\_\allowbreak{}VERSION\_\allowbreak{}MISMATCH} y cierra la conexión.
      \item El SJM muestra la \texttt{GUI\_\allowbreak{}MessageError} indicando que debe actualizarse.
      \item Fin del caso de uso.
    \end{cuenum}} \\
   & \cuCelda{\textbf{\mbox{EX-03}: Fallo al escribir en la base de datos}\par
    \begin{cuenum}[start=1]
      \item El Servidor de partidas revierte la transacción completa, de modo que no quede un \textit{Usuario} sin la \textit{Sesion} que permite usarlo (\mbox{RN-08}).
      \item El Servidor de partidas registra la excepción en la bitácora del servidor con un identificador de incidencia.
      \item El Servidor de partidas responde con \texttt{SERVER\_\allowbreak{}ERROR} y el identificador, y cierra la conexión.
      \item El SJM muestra la \texttt{GUI\_\allowbreak{}MessageError} con el identificador.
      \item El flujo continúa en el paso 1 del FN.
    \end{cuenum}} \\
   & \cuCelda{\textbf{\mbox{EX-04}: Fallo al confirmar la transacción}\par
    \begin{cuenum}[start=1]
      \item El Servidor de partidas no obtiene confirmación de que la transacción se haya escrito y la trata como fallida.
      \item El Servidor de partidas registra la excepción señalando que el estado del alta es indeterminado y responde con \texttt{SERVER\_\allowbreak{}ERROR}.
      \item El SJM muestra la \texttt{GUI\_\allowbreak{}MessageError} indicando que vuelva a intentarlo. Si el alta sí se completó, el token no llegó al cliente y esa cuenta de invitado quedará huérfana hasta caducar por el \mbox{RN-09}.
      \item Fin del caso de uso.
    \end{cuenum}} \\
   & \cuCelda{\textbf{\mbox{EX-05}: Se agota el tiempo de espera de la respuesta}\par
    \begin{cuenum}[start=1]
      \item El SJM detecta que transcurrió el tiempo límite sin recibir un mensaje completo.
      \item El SJM cierra la conexión TCP y descarta el intercambio, sin suponer que el servidor no lo procesó.
      \item El SJM registra el fallo en la bitácora local y lo informa mediante la \texttt{GUI\_\allowbreak{}MessageError}.
      \item El flujo continúa en el paso 1 del FN.
    \end{cuenum}} \\
   & \cuCelda{\textbf{\mbox{EX-06}: El mensaje recibido está malformado}\par
    \begin{cuenum}[start=1]
      \item El SJM no logra delimitar o deserializar el mensaje conforme al protocolo acordado.
      \item El SJM descarta el mensaje, cierra la conexión y registra en la bitácora local el contenido truncado.
      \item El SJM muestra la \texttt{GUI\_\allowbreak{}MessageError} indicando que la respuesta no pudo interpretarse.
      \item Fin del caso de uso.
    \end{cuenum}} \\ \hline
  \textbf{Postcondiciones} & \cuCelda{\begin{cureglas}
      \item \textbf{\mbox{POST-01}} Si el caso termina por el FN, existe un \textit{Usuario} con \texttt{tipo\_\allowbreak{}cuenta} en \texttt{INVITADO} y \texttt{estado\_\allowbreak{}cuenta} en \texttt{ACTIVA}, sin \texttt{correo} ni contraseña.
      \item \textbf{\mbox{POST-02}} Si el caso termina por el FN, ese \textit{Usuario} no tiene ninguna \textit{EstadisticaModo}: no figura en el ranking (\mbox{RN-06}).
      \item \textbf{\mbox{POST-03}} Si el caso termina por el FN, existe una \textit{Sesion} abierta y el cliente guardó su token en el dispositivo.
      \item \textbf{\mbox{POST-04}} Si el caso termina por el \mbox{FA-01}, no se creó ningún \textit{Usuario} nuevo: se abrió una \textit{Sesion} sobre la cuenta de invitado que ya existía, con su historial.
      \item \textbf{\mbox{POST-05}} Si el caso termina por cualquier excepción, no existe ningún \textit{Usuario} a medio crear.
      \item \textbf{\mbox{POST-06}} El \texttt{nickname} provisional queda ocupado y ningún otro \textit{Usuario} puede tomarlo mientras la cuenta exista.
    \end{cureglas}} \\ \hline
  \textbf{Reglas de negocio} & \cuCelda{\begin{cureglas}
      \item \textbf{\mbox{RN-01}} Antes de crear la cuenta de invitado, el SJM advierte que vive solo en este dispositivo y que sus partidas no cuentan para el ranking. En el menú principal la cuenta se distingue siempre como de invitado.
      \item \textbf{\mbox{RN-02}} El \texttt{nickname} provisional sigue un formato reservado que ningún \textit{Usuario} registrado puede elegir, de modo que una cuenta de invitado se distinga a simple vista.
      \item \textbf{\mbox{RN-03}} Una cuenta de invitado no tiene \texttt{correo} ni contraseña y, por tanto, no puede iniciar sesión con el \mbox{CU-01}. Vive atada al token guardado en el dispositivo.
      \item \textbf{\mbox{RN-04}} El token de invitado persiste en el dispositivo entre ejecuciones. Perderlo equivale a perder la cuenta: es exactamente el riesgo del que advierte el \mbox{RN-01}.
      \item \textbf{\mbox{RN-05}} Una cuenta de invitado solo entra a partidas por código: crea una sala (\mbox{CU-10}) o se une con su código (\mbox{CU-11}). Dentro de la sala y de la partida juega y chatea igual que una registrada, y como anfitrión puede expulsar (\mbox{CU-12}). Fuera de la partida puede consultar su historial desde su perfil (\mbox{CU-26}) y el perfil de otros jugadores (\mbox{CU-08}). No puede buscar partida clasificatoria (\mbox{CU-09}), consultar el ranking (\mbox{CU-25}), tener amigos (\mbox{CU-21}), invitar ni recibir invitaciones (\mbox{CU-23}), ni reportar (\mbox{CU-27}).
      \item \textbf{\mbox{RN-06}} Ninguna partida de un invitado otorga elo ni suma a \textit{EstadisticaModo}: no entra al emparejamiento, y las partidas de sala no son clasificatorias. Por eso la cuenta de invitado no tiene \textit{EstadisticaModo} y no aparece en el ranking.
      \item \textbf{\mbox{RN-07}} El número de cuentas de invitado que pueden crearse desde una misma dirección de red en una hora está acotado, para que el modo invitado no sirva para fabricar cuentas en serie.
    \end{cureglas}} \\
   & \cuCelda{\begin{cureglas}
      \item \textbf{\mbox{RN-08}} El alta del \textit{Usuario} y su \textit{Sesion} ocurren en una sola transacción: o existen las dos o no existe ninguna.
      \item \textbf{\mbox{RN-09}} Una cuenta de invitado sin actividad durante un periodo prolongado se marca como caducada, para que las cuentas huérfanas no crezcan sin límite.
      \item \textbf{\mbox{RN-10}} Una cuenta de invitado recibe strikes como cualquier otra (\mbox{RN-11} del \mbox{CU-27}), y su suspensión se levanta igual que la de una registrada: el \texttt{estado\_\allowbreak{}cuenta} vuelve a \texttt{ACTIVA} la primera vez que la cuenta se recupera después de la \texttt{fecha\_\allowbreak{}fin} de la \textit{Sancion} (\mbox{FA-06}). Mientras la sanción está vigente, el SJM conserva el token, de modo que no pueda esquivarse creando otra cuenta de invitado desde el mismo dispositivo.
    \end{cureglas}} \\ \hline
  \textbf{Observación} & \cuCelda{\textit{Sobre por qué el invitado es un Usuario y no otra cosa.}\par
    La alternativa era que el invitado existiera solo en la memoria del Servidor de partidas, sin fila en la Base de datos. Se descartó porque el invitado ocupa plazas de sala, juega partidas y escribe en el chat, y todo eso se registra en tablas que apuntan a un \textit{Usuario}. Haciendo del invitado un \textit{Usuario} desde el primer momento, los casos de sala, partida y chat funcionan igual para las dos clases de cuenta. El costo es que hay filas de \textit{Usuario} que nunca llegarán a registrarse, y por eso existe el \mbox{RN-09}.} \\ \hline
  \textbf{Datos que toca} & \cuCelda{\begin{cudatos}
      \item \textbf{\textit{Usuario}:} lee \texttt{nickname}, para comprobar que el provisional esté libre \textit{(paso 7)}
      \item \textbf{\textit{Usuario}:} crea con \texttt{tipo\_\allowbreak{}cuenta} en \texttt{INVITADO} \textit{(paso 9)}; lee el asociado al token guardado \textit{(\mbox{FA-01})}; escribe \texttt{estado\_\allowbreak{}cuenta} \textit{(\mbox{FA-06})}
      \item \textbf{\textit{Sancion}:} lee las activas de ámbito \texttt{CUENTA} \textit{(\mbox{FA-05}, \mbox{FA-06})}
      \item \textbf{\textit{Sesion}:} crea \textit{(paso 10, \mbox{FA-01})}
      \item \textbf{\textit{BitacoraAcceso}:} inserta, y lee para el límite del \mbox{RN-07} \textit{(pasos 6, 12, \mbox{FA-06})}
    \end{cudatos}} \\ \hline
\end{fichacu}
\clearpage
